\chapter{Timing and Real-Time Analysis}

\section{Introduction}

This chapter analyzes the temporal behavior of the firmware execution model. 
Although the system does not implement a real-time operating system (RTOS), 
its architecture follows deterministic scheduling principles suitable for embedded control applications.

The objective of this chapter is to:

\begin{itemize}
	\item Characterize the main control loop frequency
	\item Estimate worst-case execution time (WCET)
	\item Analyze communication latency
	\item Evaluate actuator update timing
	\item Validate system responsiveness
\end{itemize}

% ------------------------------------------------

\section{Execution Model Overview}

The firmware operates under a single-threaded cooperative execution model:

\begin{itemize}
	\item All modules execute sequentially
	\item No blocking delays are used
	\item No dynamic task switching occurs
\end{itemize}

The control loop repeats continuously with execution time determined by the cumulative duration of all module operations.

% ------------------------------------------------

\section{Control Loop Structure}

The logical execution order inside the main loop is:

\begin{enumerate}
	\item Communication processing
	\item Command decoding
	\item Control state update
	\item Motor update
	\item Servo update
	\item Sensor acquisition
	\item Telemetry packet generation
\end{enumerate}

This order ensures that new control inputs affect actuators within the same iteration cycle.

% ------------------------------------------------

\section{Worst-Case Execution Time (WCET)}

Worst-case execution time is defined as the maximum duration required to complete one full iteration of the main control loop.

Let:

\[
T_{loop} = T_{comm} + T_{control} + T_{motor} + T_{servo} + T_{sensor} + T_{telemetry}
\]

Each term represents the maximum processing time of the corresponding subsystem.

\subsection{Qualitative Estimation}

In typical operation:

\begin{itemize}
	\item Communication parsing is lightweight (byte buffer inspection)
	\item Motor and servo updates involve simple register writes
	\item Sensor acquisition depends on I2C latency
	\item Telemetry formatting involves structured serialization
\end{itemize}

The dominant time contributor is typically I2C communication when accessing external peripherals.

% ------------------------------------------------

\section{Loop Frequency Estimation}

The control loop frequency is approximated as:

\[
f_{loop} = \frac{1}{T_{loop}}
\]

For responsive RC control systems, a target loop frequency above 50 Hz is sufficient. 
Higher frequencies (100–500 Hz) improve control smoothness and reduce perceptible latency.

Given the lightweight nature of operations and the processing capability of the microcontroller, 
the system is expected to operate significantly above the minimum 50 Hz requirement.

% ------------------------------------------------

\section{Communication Latency}

End-to-end command latency consists of:

\begin{itemize}
	\item Transmission delay over Bluetooth
	\item Packet parsing delay
	\item Control loop synchronization delay
\end{itemize}

Total latency can be approximated as:

\[
T_{latency} = T_{bt} + T_{parse} + T_{loop\_sync}
\]

Where:

\begin{itemize}
	\item $T_{bt}$ is wireless transmission delay
	\item $T_{parse}$ is packet validation time
	\item $T_{loop\_sync}$ is the maximum time until next loop iteration
\end{itemize}

Because the loop runs continuously without blocking delays, 
$T_{loop\_sync}$ is bounded by $T_{loop}$.

This guarantees predictable maximum response delay.

% ------------------------------------------------

\section{Actuator Update Timing}

Motor and servo updates occur once per loop iteration.

Therefore:

\[
T_{actuator\_update} \leq T_{loop}
\]

This ensures:

\begin{itemize}
	\item Deterministic actuator refresh
	\item No irregular pulse timing
	\item Stable PWM output generation
\end{itemize}

Servo control timing is hardware-assisted by internal PWM peripherals, 
reducing dependency on software timing precision.

% ------------------------------------------------

\section{Sensor Sampling Rate}

Sensor sampling frequency is also determined by loop frequency:

\[
f_{sensor} = f_{loop}
\]

However, for sensors requiring longer acquisition time (e.g., via I2C), 
effective sampling frequency may be slightly lower.

Sensor reads are designed to:

\begin{itemize}
	\item Avoid blocking delays
	\item Complete within a bounded execution time
\end{itemize}

% ------------------------------------------------

\section{Determinism Analysis}

The firmware achieves determinism through:

\begin{itemize}
	\item Fixed execution order
	\item Absence of random delays
	\item No dynamic memory allocation
	\item Bounded buffer sizes
\end{itemize}

This ensures repeatable timing behavior under identical operating conditions.

% ------------------------------------------------

\section{Jitter Considerations}

Jitter refers to variations in loop execution time.

Primary jitter sources include:

\begin{itemize}
	\item Variable Bluetooth packet arrival timing
	\item I2C bus contention
	\item Interrupt-driven background processes
\end{itemize}

However, because:

\begin{itemize}
	\item Actuator signals use hardware PWM
	\item No blocking communication is used
\end{itemize}

Control stability is not significantly affected by small loop timing variations.

% ------------------------------------------------

\section{Real-Time Classification}

The system can be classified as:

\begin{itemize}
	\item Soft real-time
\end{itemize}

Reasoning:

\begin{itemize}
	\item Occasional minor delays do not cause catastrophic failure
	\item Deadlines are important for performance but not safety-critical
	\item System maintains stable state during communication loss
\end{itemize}

This classification is appropriate for RC control systems.

% ------------------------------------------------

\section{Communication Loss Handling}

If communication is interrupted:

\begin{itemize}
	\item Last valid control state is retained
	\item System remains stable
	\item Optional timeout-based fail-safe may be implemented
\end{itemize}

Timeout detection can be expressed as:

\[
T_{now} - T_{last\_packet} > T_{timeout}
\]

If true, actuators can be set to predefined safe states.

% ------------------------------------------------

\section{Scalability Impact on Timing}

Adding new modules increases:

\[
T_{loop}
\]

However, because of modular structure:

\begin{itemize}
	\item New features can be evaluated individually
	\item Timing impact can be measured incrementally
	\item System limits remain predictable
\end{itemize}

This supports controlled expansion without architectural degradation.

% ------------------------------------------------

\section{Summary}

The firmware timing model provides:

\begin{itemize}
	\item Deterministic execution order
	\item Bounded worst-case loop duration
	\item Predictable actuator update timing
	\item Controlled communication latency
	\item Soft real-time operational classification
\end{itemize}

The next chapter presents validation testing and performance measurements to empirically verify these timing characteristics.